\documentclass[10pt,a4paper]{amsart}
\usepackage[margin=19mm]{geometry}
\usepackage{amsmath,amssymb,mathtools}
\usepackage[scr=rsfs,cal=euler]{mathalfa}
\usepackage{xfrac}
\usepackage[unicode,hidelinks,bookmarks=false]{hyperref}
\newcommand{\ie}{i.e.\ }
\newcommand{\cf}{cf.\ }
\newcommand{\resp}{respectively\ }
\newcommand{\Subsection}{\subsection}
\providecommand{\nobreakdash}{\nobreak\hbox{-}}
\setlength{\emergencystretch}{2em}
\multlinegap0pt
\usepackage{amssymb}
\usepackage{xcolor}
\usepackage[shortlabels]{enumitem}
\usepackage{tikz}
\newtheorem{proposition}{Proposition}
\def\e{\mathrm{e}}
\title{Inspections}
\author{Refael Hassin, Liron Ravner}
\date{Source: arXiv:2606.02464v1 (CC BY 4.0)}
\begin{document}
\maketitle

\noindent\emph{Source and licence.} Inspections. Version arXiv:2606.02464v1 (CC BY 4.0), distributed by arXiv under the Creative Commons Attribution 4.0 International licence, http://creativecommons.org/licenses/by/4.0/, as recorded in the arXiv licence metadata for this exact version and shown on the abstract page https://arxiv.org/abs/2606.02464v1. Full licence notice: \texttt{licenses/arxiv-2606.02464-CC-BY-4.0.txt}.\par\smallskip

\noindent\textit{Editorial scope.} Counted unit: the article's proposition prop:rho\_dominant (source0.tex lines 132--137). The article's complete original proof of it is delivered with it (source0.tex lines 138--179, one \texttt{proof} environment of 42 lines). No mathematics is written, repaired or re-derived: the statement and the proof are the article's own lines, in the article's own notation, and no retained line is substituted or re-flowed.  No further source-local context is needed: every cross-reference of every retained line resolves inside the two delivered spans.  Nothing was omitted from any retained span, so the package carries no omission record.  No retained line contains a citation, so the wrapper carries no bibliography and no bibliography entry.  No retained line contains a figure environment, an image-inclusion command or drawing code, so no image is delivered and the package declares no asset dependency.  Editorial formatting only, in addition: an AMS \texttt{amsart} preamble that restates the article's own declarations and macros used by the retained lines, namely the environments \texttt{proposition} on separate counters, together with the standard packages those lines need.  The retained environments print in this document as Proposition 1 in order of appearance, whereas the article prints the same items with the numbering of its own full document.  The complete native source of the article as distributed in the arXiv e-print for this exact version is bundled unchanged as \texttt{sources/201904/source0.tex}.
\begin{proposition}\label{prop:rho_dominant}
 If $\lambda_\tau=\lambda p$ and $p\in(0,1)$, then
 \begin{equation}\label{eq:rho_dominant}
      \rho>\sigma\ .
\end{equation}
\end{proposition}

\begin{proof}
%to establish that for any RR policy $p\in[0,1]$ there exists a GA policy with a smaller $\sigma$ $\tau$ achieving a higher utility. 
Fix $\lambda,\mu>0$ and let $p\in(0,1)$ be such that $\rho(p)=\lambda p/\mu<1$. The GA threshold $\tau> 0$ is chosen so that the arrival rates are equal under both policies;
\begin{equation}\label{eq:equal_rates}
    \lambda p=\lambda_\tau=\frac{\lambda}{1+\lambda\tau}\quad\Longrightarrow\quad
p=\frac{1}{1+\lambda\tau}\ .
\end{equation}
Let $\tilde a(s)$ be the LST of an exponential random variable with rate $\lambda p$. Then, by \eqref{eq:equal_rates},
\begin{align*}
\tilde a(s)=\frac{\lambda p}{\lambda p+s}=\frac{\lambda}{\lambda+s(1+\lambda\tau)} \ .
\end{align*}
For any $s\geq 0$, we have that
\begin{equation}\label{eq:a_inequality}
    a(s)=\frac{\lambda e^{-s\tau}}{\lambda+s}
\leq \frac{\lambda}{(\lambda+s)(1+\tau s)}=\frac{\lambda}{\lambda+s+\lambda s\tau+\tau s^2}
\leq \frac{\lambda}{\lambda+s(1+\lambda\tau)}= \tilde a(s)\ ,
\end{equation}
where the exponential inequality $\e^{-x}\leq \frac{1}{1+x}$ for all $x\geq 0$ was used in the first inequality. Note that the second inequality becomes strict if $s>0$.

As M/M/1 is a special case of G/M/1, $\rho$ is also a solution of a fixed point equation. That is,
\begin{align*}
    \rho=\tilde{a}(\mu(1-\rho))\ ,
\end{align*}
and 
\begin{align*}
    \sigma=a(\mu(1-\sigma))\ .
\end{align*}
Let $g(x):=a(\mu(1-x))$. From the properties of the LST, $g$ is a convex function that increases with $x>0$ and satisfies
\begin{align*}
    g(0)=a(\mu)\in(0,1)\  , g(1)=a(0)=1\ .
\end{align*}
The same is true for the function $f(x):=\tilde{a}(\mu(1-x))$ with the boundary values
\begin{align*}
    f(0)=\tilde{a}(\mu)\geq a(\mu)\  , f(1)=\tilde{a}(0)=1\ ,
\end{align*}
where the inequality follows from \eqref{eq:a_inequality}.
Hence, both functions admit a unique fixed point in $x\in(0,1)$. Suppose that $\sigma=g(\sigma)$, then by \eqref{eq:a_inequality},
\begin{align*}
    \sigma=g(\sigma)<f(\sigma)\ ,
\end{align*}
which implies that the solution of $\rho=f(\rho)$ has to be larger than $\sigma$, i.e., $\rho>\sigma$. 
\end{proof}

\end{document}
